\subsection{Несобственный интеграл Римана}

\begin{definition}
	Пусть $f$ интегрируема по Риману на $[a; \tilde{b}],\ \forall \tilde{b} \in [a; b)$. Тогда интеграл $\int_a^b f(x)dx$ называется \textit{несобственным интегралом Римана}
	\begin{itemize}
		\item \textit{первого рода}, если $b = +\infty$
		
		\item \textit{второго рода}, если $b < +\infty$ и $f$ не является ограниченной на $[a; b)$
	\end{itemize}
	Его понимают как $\liml_{\tilde{b} \to b-0} \int_a^{\tilde{b}} f(x)dx$ или $\liml_{\tilde{b} \to +\infty} \int_a^{\tilde{b}} f(x)dx$. Если этот предел существует и конечен, то несобственный интеграл Римана соответствующего рода \textit{сходится} и этот предел равен его значению. В противном случае, интеграл \textit{расходится}.
	\\
	Аналогично можно определить интеграл с невключённым левым и включённым правым концом интегрирования.
\end{definition}

\begin{anote}
	То есть, когда мы написали интеграл, мы ещё не знаем скрывается за ним какое-то значение и он  означает число, или же он расходится, но написать его мы можем и без этого.
\end{anote}

\begin{definition}
	Если при стремлении к бесконечности функция неограниченна, то \textit{главное значение в смысле Коши} $v.p.\int_{-\infty}^{+\infty} f(x)dx = \liml_{A\to+\infty}\int_{-A}^A f(x)dx$, аналогично,если $f$ определена на  $(b,c)$, и на нём есть точка $a$ в окрестности которой функция неограниченна. Тогда $v.p.\int_{b}^{c} f(x)dx = \liml_{\eps\to+0}\int_{b}^{a-\eps} f(x)dx + \int_{a + \eps}^{c} f(x)dx$. Мы как бы пытаемся посчитать всё без этой точки.
\end{definition}